\section{Przegląd literatury i~analiza istniejących rozwiązań}
\label{chap:pierwszy}

Rozdział przedstawia przegląd literatury w~obszarach kluczowych dla systemu RiskScoop AI: agentów AI w~analizie geoprzestrzennej, systemów monitorowania powodzi oraz otwartych danych geoprzestrzennych.

\subsection{Agenci sztucznej inteligencji w~analizie geoprzestrzennej}

\subsubsection{Koncepcja agentów AI}

Agenty sztucznej inteligencji definiuje się jako autonomiczne systemy zdolne do percepcji otoczenia, podejmowania decyzji oraz wykonywania działań~\cite{google2024geospatial}. W~kontekście analizy geoprzestrzennej, agenty AI pośredniczą między użytkownikiem formułującym zapytanie w~języku naturalnym a~złożonymi operacjami przetwarzania danych przestrzennych.

Architektura współczesnych agentów geoprzestrzennych opiera się na trzech komponentach:
\begin{itemize}
    \item \textbf{Moduł percepcji} -- rozumienie języka naturalnego i~ekstrakcja parametrów
    \item \textbf{Moduł rozumowania} -- planowanie sekwencji operacji geoprzestrzennych
    \item \textbf{Moduł działania} -- wykonanie operacji i~generowanie wyników
\end{itemize}

\subsubsection{Przegląd istniejących frameworków}

W~literaturze opisano kilka frameworków agentów AI dla analiz geoprzestrzennych. \textbf{ChatGeoAI}~\cite{li2024chatgeoai} wykorzystuje model Llama~2 do generowania kodu PyQGIS, osiągając wskaźnik sukcesu 76\%. \textbf{LLM-Find}~\cite{gong2024llmfind} specjalizuje się w~autonomicznym wyszukiwaniu danych geoprzestrzennych z~80--90\% skutecznością. \textbf{GeoGPT}~\cite{rao2024geogpt} oferuje kompleksowe podejście end-to-end do analiz przestrzennych. Najnowszy \textbf{GeoAgent}~\cite{li2024geoagent} integruje Monte Carlo Tree Search z~RAG, osiągając 87\% sukcesu.

Tabela~\ref{tab:ai-agents-comparison} przedstawia porównanie frameworków.

\begin{table}[H]
    \centering
    \caption{Porównanie frameworków agentów AI w~analizie geoprzestrzennej}
    \begin{tabular}{|l|c|l|}
        \hline
        \textbf{Framework} & \textbf{Sukces} & \textbf{Innowacja} \\
        \hline
        ChatGeoAI & 76\% & NER + PyQGIS \\
        \hline
        LLM-Find & 80--90\% & Autonomiczne debugowanie \\
        \hline
        GeoAgent & 87\% & MCTS + RAG \\
        \hline
    \end{tabular}
    \label{tab:ai-agents-comparison}
\end{table}

Żaden z~frameworków nie specjalizuje się w~analizie ryzyka powodziowego dla infrastruktury krytycznej.

\subsubsection{Framework Agno}

Do implementacji RiskScoop AI wybrano framework Agno~\cite{agno2024} ze względu na:
\begin{itemize}
    \item Wydajność -- instancjonowanie agenta w~3~$\mu$s (529× szybciej niż LangGraph)
    \item Niskie zużycie pamięci -- 6.6~KiB na agenta
    \item Wbudowaną integrację z~FastAPI (AgentOS)
    \item Natywne wsparcie dla modeli Google Gemini
\end{itemize}

\subsection{Systemy monitorowania powodzi}

\subsubsection{Copernicus GloFAS}

Global Flood Awareness System (GloFAS) to operacyjny system monitorowania powodzi prowadzony przez Komisję Europejską~\cite{jrc2020glofas}. System bazuje na modelu hydrologicznym LISFLOOD oraz danych meteorologicznych ERA5, generując prognozy przepływu rzek na okres do 30 dni.

GloFAS oferuje mapy zagrożenia powodziowego dla różnych okresów powrotnych (RP10, RP20, RP50, RP100), gdzie okres powrotny oznacza statystyczny czas między zdarzeniami o~określonej skali. Dane klasyfikowane są w~5 klasach głębokości zalania (tab.~\ref{tab:glofas_classes}).

\begin{table}[H]
    \centering
    \caption{Klasy głębokości GloFAS~\cite{glofas2024}}
    \begin{tabular}{|c|l|l|}
        \hline
        \textbf{Klasa} & \textbf{Głębokość} & \textbf{Charakterystyka} \\
        \hline
        1 & 0--0,5~m & Niskie zalanie \\
        \hline
        2 & 0,5--1~m & Zagrożenie dla pieszych \\
        \hline
        3 & 1--2~m & Zagrożenie dla budynków \\
        \hline
        4 & 2--5~m & Poważne zniszczenia \\
        \hline
        5 & >5~m & Katastrofalne zalanie \\
        \hline
    \end{tabular}
    \label{tab:glofas_classes}
\end{table}

\subsubsection{Copernicus Rapid Mapping}

Copernicus Emergency Management Service oferuje usługę Rapid Mapping dostarczającą mapy sytuacyjne w~ciągu kilku godzin od zgłoszenia~\cite{copernicus2024ems}. Średni czas dostawy pierwszej mapy wynosi 4,2 godziny. Usługa jest aktywowana na wniosek służb kryzysowych i~wykorzystuje obrazy satelitarne wysokiej rozdzielczości.

\subsection{Otwarte dane geoprzestrzenne}

\subsubsection{Overture Maps Foundation}

Overture Maps Foundation~\cite{overture2024foundation} to inicjatywa Linux Foundation wspierana przez Amazon, Meta, Microsoft i~TomTom, udostępniająca otwarte dane geoprzestrzenne:
\begin{itemize}
    \item \textbf{Buildings} -- 2,6 miliarda budynków z~atrybutami (wysokość, funkcja)
    \item \textbf{Places} -- 64 miliony POI z~kategoryzacją (szpitale, szkoły, transport)
    \item \textbf{Transportation} -- 321 milionów segmentów drogowych
    \item \textbf{Admins} -- granice administracyjne wszystkich krajów
\end{itemize}

Dane są dystrybuowane w~formacie GeoParquet i~dostępne przez dedykowane API.

\subsubsection{OpenStreetMap i~Nominatim}

OpenStreetMap (OSM) to crowdsourcingowy projekt mapowy dostarczający danych o~granicach administracyjnych. Nominatim to usługa geokodowania OSM umożliwiająca wyszukiwanie lokalizacji po nazwie. Overpass API pozwala na pobieranie szczegółowych geometrii granic.

\subsection{Model oceny ryzyka powodziowego}

Tradycyjny model oceny ryzyka katastroficznego~\cite{pregnolato2022bridge} definiuje ryzyko jako funkcję trzech komponentów: zagrożenia (hazard), ekspozycji (exposure) i~podatności (vulnerability). W~praktyce pełna implementacja takiego modelu wymaga szczegółowych danych o~podatności poszczególnych typów obiektów, które często nie są dostępne.

RiskScoop AI stosuje uproszczone podejście oparte na danych Copernicus GloFAS: klasy głębokości zalania (1--5) z~map zagrożenia powodziowego są mapowane na trójstopniową skalę ryzyka (niskie, średnie, wysokie), a~ekspozycja jest wyznaczana poprzez analizę przecięć przestrzennych obiektów infrastruktury ze strefami zagrożenia. Podejście to pomija komponent podatności, lecz pozwala na szybką identyfikację zagrożonych obiektów w~kontekście Rapid Mapping.

W~kontekście infrastruktury krytycznej~\cite{chukwuma2021flood} szczególnie istotne są: służba zdrowia (szpitale, kliniki), transport (drogi, mosty), edukacja (szkoły) oraz usługi publiczne (straż pożarna, policja).

\subsection{Podsumowanie przeglądu literatury}

Z~przeglądu literatury wynika, że nie istnieje system łączący interfejs języka naturalnego z~analizą ryzyka powodziowego dla infrastruktury krytycznej. Istniejące frameworki agentów AI są generalistyczne, a~platformy flood risk nie oferują interfejsu NL. RiskScoop AI wypełnia tę lukę, specjalizując się w~Rapid Mapping z~wykorzystaniem otwartych europejskich źródeł danych.

\subsection{Analiza istniejących rozwiązań}

Podrozdział przedstawia analizę konkurencyjnych rozwiązań na rynku systemów analizy geoprzestrzennej z~interfejsem języka naturalnego oraz platform oceny ryzyka powodziowego.

\subsubsection{Platformy geoinformacyjne z~interfejsem języka naturalnego}

\paragraph{Ageospatial}

Ageospatial to szwajcarski startup oferujący platformę analizy geoprzestrzennej z~agentem AI, rozwinięty we współpracy ze Swiss Federal Office of Topography (swisstopo). Platforma obsługuje zapytania typu ,,Budynki zagrożone osuwiskami'' czy ,,Ryzyko powodziowe w~okolicy mojego domu''. Oferuje wdrożenia lokalne (on-premise) dla instytucji rządowych oraz visible reasoning graphs zwiększające przejrzystość decyzji agenta.

\paragraph{Ask.Earth}

Ask.Earth oferuje platformę GIA (Geospatial Intelligence Agent) przetwarzającą zapytania w~języku naturalnym dotyczące obrazów satelitarnych. Platforma koncentruje się na sektorach: energetyka (współpraca z~TotalEnergies), monitoring środowiskowy i~intelligence przedsiębiorstw. Nie specjalizuje się w~ryzyku powodziowym ani infrastrukturze krytycznej.

\paragraph{Google Earth AI z~Gemini}

W~2024 roku Google uruchomił pilotaż Google Earth AI z~modelem Gemini, oferując geospatial reasoning oraz integrację z~Google Earth Engine. Wejście Google stanowi zagrożenie dla platform generalistycznych, jednak specjalizacja niszowa (flood risk dla infrastruktury) pozostaje poza zainteresowaniem giganta technologicznego.

\subsubsection{Platformy oceny ryzyka powodziowego}

\paragraph{SaferPlaces}

SaferPlaces oferuje Digital Twin dla analizy ryzyka powodziowego, łącząc szybkie algorytmy AI z~modelami fizycznymi. System generuje wyniki w~czasie rzeczywistym dla obszarów od pojedynczych posesji po regiony metropolitalne. Skierowany jest do ubezpieczycieli i~administracji publicznej. Nie oferuje interfejsu języka naturalnego.

\paragraph{Blue Sky Analytics}

Blue Sky Analytics specjalizuje się w~monitorowaniu powodzi z~danych satelitarnych, transformując terabajty danych w~actionable datasets. Koncentruje się na near real-time monitoring aktualnych powodzi, podczas gdy RiskScoop AI oferuje pre-disaster risk assessment.

\paragraph{Flood Factor (First Street Foundation)}

Flood Factor to narzędzie non-profit oferujące ocenę ryzyka powodziowego dla 142 milionów nieruchomości w~USA~\cite{firststreet2020flood}. Przypisuje każdej nieruchomości Flood Factor Score (1--10) uwzględniając projekcje klimatyczne na 30 lat. Pokrycie geograficzne ograniczone do USA, brak interfejsu NL, koncentracja na nieruchomościach mieszkalnych.

\subsubsection{Prototypy akademickie}

Prototypy akademickie (ChatGeoAI, LLM-Find, GeoGPT, GeoAgent) opisane wcześniej w~niniejszym rozdziale demonstrują techniczną wykonalność integracji LLM z~operacjami geoprzestrzennymi. Osiągają wskaźniki sukcesu 76--87\%, co stanowi benchmark dla RiskScoop AI. Żaden nie specjalizuje się w~flood risk assessment.

\subsubsection{Analiza porównawcza}

Tabela~\ref{tab:competitive-comparison} przedstawia porównanie kluczowych rozwiązań.

\begin{table}[H]
\centering
\caption{Porównanie konkurencyjnych rozwiązań}
\label{tab:competitive-comparison}
\begin{tabular}{|l|c|c|c|c|}
\hline
\textbf{Rozwiązanie} & \textbf{NL} & \textbf{Flood} & \textbf{Infra} & \textbf{Typ} \\
\hline
Ageospatial & Tak & Częściowo & Nie & Komercyjny \\
\hline
Ask.Earth & Tak & Nie & Nie & Komercyjny \\
\hline
Google Earth AI & Tak & Nie & Nie & Korporacyjny \\
\hline
SaferPlaces & Nie & Tak & Częściowo & Komercyjny \\
\hline
Blue Sky Analytics & Nie & Tak & Nie & Komercyjny \\
\hline
Flood Factor & Nie & Tak (USA) & Nie & Non-profit \\
\hline
ChatGeoAI & Tak & Nie & Nie & Akademicki \\
\hline
GeoAgent & Tak & Nie & Nie & Akademicki \\
\hline
\textbf{RiskScoop AI} & \textbf{Tak} & \textbf{Tak} & \textbf{Tak} & \textbf{Akademicki} \\
\hline
\end{tabular}
\end{table}

Legenda: NL -- interfejs języka naturalnego, Flood -- ocena ryzyka powodziowego, Infra -- specjalizacja w~infrastrukturze krytycznej.

\subsubsection{Identyfikacja luki rynkowej}

Z~analizy wynika, że \textbf{żadne istniejące rozwiązanie nie łączy interfejsu języka naturalnego z~oceną ryzyka powodziowego dla infrastruktury krytycznej}. Platformy NL (Ageospatial, Ask.Earth) nie specjalizują się w~flood risk. Platformy flood risk (SaferPlaces, Flood Factor) nie oferują interfejsu NL. Żadne rozwiązanie nie koncentruje się na infrastrukturze krytycznej w~kontekście Rapid Mapping.

RiskScoop AI wypełnia tę lukę poprzez:
\begin{itemize}
    \item Interfejs języka naturalnego dla intuicyjnych zapytań
    \item Integrację z~Copernicus GloFAS dla danych powodziowych
    \item Specjalizację w~infrastrukturze krytycznej (szpitale, szkoły, transport)
    \item Wykorzystanie europejskich otwartych danych (Overture Maps, GloFAS)
\end{itemize}

\subsection{Wymagania funkcjonalne i~niefunkcjonalne}

System RiskScoop AI realizuje następujący zakres funkcjonalny:

\paragraph{Przetwarzanie języka naturalnego}
System przyjmuje zapytania użytkownika w~języku naturalnym dotyczące konkretnych lokalizacji oraz typów obiektów infrastrukturalnych (szpitale, drogi, mosty, szkoły, obiekty użyteczności publicznej). Agent AI automatycznie tłumaczy zapytania na operacje geoprzestrzenne, identyfikując:
\begin{itemize}
    \item typ infrastruktury (kategoria obiektu),
    \item obszar geograficzny (miasto, województwo, kraj),
    \item intencję użytkownika (identyfikacja zagrożonych obiektów, kwantyfikacja ryzyka).
\end{itemize}

Przykładowe zapytania obsługiwane przez system:
\begin{itemize}
    \item ,,Które szpitale w~Warszawie są zagrożone powodzią?'',
    \item ,,Pokaż drogi krajowe w~strefie wysokiego ryzyka powodziowego w~województwie małopolskim'',
    \item ,,Ile szkół podstawowych znajduje się w~obszarze zagrożenia powodziowego w~Krakowie?''.
\end{itemize}

\paragraph{Integracja źródeł danych}
\begin{itemize}
    \item \textbf{Overture Maps:} Baza danych infrastrukturalnych zawierająca informacje o~budynkach, punktach użyteczności publicznej (POI), sieci transportowej oraz granicach administracyjnych. System wykorzystuje tematyczne warstwy danych (Buildings, Places, Transportation, Admins) do identyfikacji obiektów infrastruktury krytycznej.

    \item \textbf{Copernicus GloFAS:} Globalny System Świadomości Powodzi dostarczający dane o~zagrożeniu powodziowym, w~tym prognozy przepływu rzek oraz mapy zasięgu potencjalnych powodzi. System integruje dane GloFAS w~celu określenia poziomu zagrożenia dla zidentyfikowanych obiektów.
\end{itemize}

\paragraph{Kwantyfikacja ryzyka powodziowego}
System klasyfikuje ryzyko powodziowe dla zidentyfikowanych obiektów infrastrukturalnych na podstawie danych z~Copernicus GloFAS. Klasy głębokości zalania (1--5) z~map zagrożenia powodziowego są mapowane na trójstopniową skalę ryzyka: niskie, średnie, wysokie. Dla każdego obiektu system wyznacza poziom ryzyka poprzez analizę przecięć przestrzennych z~punktami powodziowymi w~zadanym buforze odległości.

\paragraph{Wizualizacja wyników}
System prezentuje wyniki analizy na interaktywnej mapie wykorzystującej bibliotekę Mapbox GL JS, oferując:
\begin{itemize}
    \item kodowanie kolorystyczne obiektów według poziomu ryzyka,
    \item warstwy danych o~zagrożeniu powodziowym (GloFAS),
    \item szczegółowe informacje o~każdym zidentyfikowanym obiekcie (nazwa, typ, adres, poziom ryzyka),
    \item funkcje filtrowania i~wyszukiwania wyników,
    \item interaktywną nawigację i~kontrolę powiększenia mapy.
\end{itemize}

\paragraph{Interfejs użytkownika}
Aplikacja webowa zapewnia responsywny interfejs dostępny przez przeglądarkę internetową, składający się z:
\begin{itemize}
    \item pola wprowadzania zapytania w~języku naturalnym,
    \item panelu wyników z~listą zidentyfikowanych obiektów,
    \item interaktywnej mapy z~wizualizacją ryzyka,
    \item panelu szczegółów obiektu z~informacjami o~poziomie zagrożenia.
\end{itemize}

